\section{引言：从 Model 到 Agent}

林俊旸（Qwen 团队负责人）在清华-北大联合实验室的活动中，介绍了通义千问（Qwen）2025 年的技术进展。演讲主题从最初的 ``Towards a Generalist Model'' 改为 ``Towards a Generalist Agent''——这一改变反映了整个 AI 领域的范式转移：从``训练更强的模型''到``构建能自主使用工具的智能体''。

\begin{importantbox}{为什么从 Model 到 Agent？}
林俊旸认为，Agent 是一个比 Model 更大的概念。人类与动物的核心差异之一在于\textbf{自主使用工具的能力}。而今天的训练范式变化（特别是强化学习的引入）使得 AI 也获得了这种可能：只要解决了推理和评估，模型就能被训练去执行各种任务——无论是 Digital Agent 还是 Physical Agent。
\end{importantbox}

\begin{knowledgebox}{Qwen 开源生态概览}
\begin{itemize}
    \item 2023年8月3日开始开源之路
    \item 体验入口：chat.qwen.ai（聚合闭源和开源模型）
    \item 模型发布平台：Hugging Face 和 ModelScope（内容同步）
    \item GitHub 以脚本为主；技术博客代替论文（如 Qwen3-Next 架构等）
    \item 吉祥物是一只水豚（Capybara），取名佛系风格
\end{itemize}
\end{knowledgebox}

\subsection{本章小结}

Qwen 团队的愿景是构建一个 ``Multimodal Foundation Agent''——一个能看（视觉）、能听（语音）、能说（生成）、能做（工具使用）的通用智能体。这一愿景从 2020 年的多模态研究延续至今。

